\section{Dos rutas técnicas: In-Context y End-to-End}

La sección anterior definió el Agent y los tipos de entorno; esta sección trata de cómo construir un Agent. El programa divide las rutas en In-Context Learning y End-to-End Training. La ruta in-context se apoya en un modelo preentrenado potente, el prompt, los protocolos de herramientas y la ingeniería de contexto; permite iterar rápido y ofrece mucha flexibilidad. La ruta end-to-end incorpora la conducta en los pesos del modelo; la inferencia es más estable, pero el costo de entrenamiento es alto y construir los entornos es complejo.

\begin{figure}[H]
\centering
\includegraphics[width=0.94\textwidth]{agent-training-routes.png}
\caption{Dos rutas de entrenamiento de un Agent: In-context es flexible; End-to-end es estable, pero de alto costo. Diagrama conceptual propio, elaborado a partir de la conversación de 00:14:50--00:30:57.}
\end{figure}

\begin{knowledgebox}{Lectura de la figura: las dos rutas no son excluyentes}
Muchos sistemas reales combinan ambas: usan in-context para montar el producto con rapidez, usan trayectorias sintéticas y RL para entrenar en el modelo las conductas clave, y después usan la ingeniería de contexto para resolver los detalles de producción.
\end{knowledgebox}

\subsection{Los tres componentes de Agent Training}

Toda ruta técnica termina en el entrenamiento y el despliegue; esta subsección desglosa los tres componentes de Agent training. Sus etapas clave son la síntesis de datos, el aprendizaje por refuerzo y la seguridad. La síntesis de datos genera trajectory de alta calidad; el aprendizaje por refuerzo depende de una task clara y de una verifiable reward; la seguridad requiere sandbox, restricciones de conducta y human-in-the-loop.

\begin{figure}[H]
\centering
\includegraphics[width=0.96\textwidth]{training-ingredients.png}
\caption{Los tres componentes de Agent Training: las trayectorias sintéticas, el aprendizaje por refuerzo y las restricciones de seguridad forman juntos el pipeline de entrenamiento. Diagrama conceptual propio, elaborado a partir de la conversación de 00:28:29--00:30:57.}
\end{figure}

\begin{importantbox}{El cambio de paradigma en los datos de un Agent}
Los datos anotados tradicionales suelen ser input-output; los datos de un Agent se parecen más a environment + task + action trajectory + reward. El modelo no solo aprende respuestas, sino cómo actuar dentro de un entorno.
\end{importantbox}

\begin{knowledgebox}{Comparación de las etapas de entrenamiento}
\begin{tabularx}{\textwidth}{p{0.22\textwidth}p{0.34\textwidth}X}
\toprule
Etapa & Producto & Riesgo principal \\
\midrule
Data Synthesis & Trayectorias de acción de alta calidad y demostraciones de llamadas a herramientas & Si las trayectorias no son realistas, el modelo aprende a actuar un papel. \\
RL & Optimizar la política a partir de la retroalimentación del entorno & La reward está mal diseñada o sufre un hack. \\
Safety & Entorno aislado, permisos, confirmación humana, políticas de rechazo & Si es demasiado laxa, es peligrosa; si es demasiado estricta, la tarea no se puede completar. \\
Evaluation & Tasa de éxito de las tareas, costo, capacidad de recuperación & El benchmark no coincide con las tareas reales. \\
\bottomrule
\end{tabularx}
\end{knowledgebox}

\subsection{Resumen de la sección}

Lo central de Agent training es colocar la tarea dentro de un entorno para que el modelo aprenda a actuar mediante trayectorias, recompensas y restricciones de seguridad. No es solo SFT ni solo prompt engineering.
